\chapter{Cơ sở lý thuyết về hệ thống số Posit và phép toán MAC}
\label{ch:arithmetic}

\section{Lý do lựa chọn hệ thống số Posit(32,3) thay cho IEEE-754}

Trong thiết kế phần cứng số học dấu phẩy động 32 bit, chuẩn
\texttt{IEEE-754 binary32 (FP32)} \cite{ieee754_2019} sử dụng cấu trúc cố định
gồm 1 bit dấu, 8 bit mũ có độ lệch ($bias=127$) và 23 bit fraction (24 bit
significand kể cả bit ẩn). Mặc dù phổ biến trên bộ xử lý đa dụng, cấu trúc
trường cố định này bộc lộ một số hạn chế khi triển khai các phép toán tích vô
hướng (\textit{dot-product} / MAC) chuyên dụng trên FPGA:
\begin{itemize}
  \item \textbf{Lãng phí không gian mã cho các trạng thái ngoại lệ:} IEEE-754 tách
        hai biểu diễn khác nhau cho số không ($+0$ và $-0$), dùng hai mã vô cực
        ($+\infty, -\infty$) và hàng triệu tổ hợp bit cho \texttt{NaN}. Ngoài ra,
        miền số dưới chuẩn (\textit{subnormal/denormal}, có bit ẩn bằng 0) đòi
        hỏi mạch phát hiện và chuẩn hóa riêng biệt, làm tăng độ phức tạp giải mã
        và đóng gói phần cứng.
  \item \textbf{Phân bố độ phân giải phẳng theo hàm mũ:} FP32 giữ cố định 23 bit
        fraction ở mọi bậc mũ bình thường, do đó không ưu tiên thêm độ chính xác
        cho miền giá trị phổ biến quanh $\pm 1$ thường gặp trong chuẩn hóa trọng
        số lọc số và mạng nơ-ron.
\end{itemize}

Hệ thống số \textbf{Posit} \cite{gustafson2017posit,positstandard2022} khắc phục
các nhược điểm trên bằng cách đưa thêm trường \textit{regime} có độ dài biến thiên
trên cùng từ mã 32 bit (lưu ý chuẩn Posit 2022 ấn định $es=2$ làm chuẩn công
nghiệp, còn đồ án sử dụng cấu hình mở rộng tùy biến \textbf{\texttt{Posit(32,3)}}
với $n=32, es=3$ theo Norris--Kim \cite{norris2021iterativeposit}). Trên cơ sở
phân tích định lượng giữa \texttt{FP32}, \texttt{Posit(32,2)} và \texttt{Posit(32,3)}
(Bảng~\ref{tab:ieee-vs-posit}), đề tài lựa chọn \texttt{Posit(32,3)} vì bốn lý do
kỹ thuật trực tiếp sau:
\begin{enumerate}
  \item \textbf{Dải động vượt trội hơn FP32:} Cần lưu ý rằng không phải mọi cấu
        hình Posit 32 bit đều có dải biểu diễn rộng hơn FP32. Với $es=2$ của
        chuẩn 2022, \texttt{Posit(32,2)} chỉ đạt $\mathrm{maxpos}=2^{120}\approx
        1{,}3292\times 10^{36}$, nhỏ hơn giá trị hữu hạn lớn nhất của FP32
        ($(2-2^{-23})\times 2^{127}\approx 3{,}4028\times 10^{38}$) khoảng $256$ lần.
        Ngược lại, khi mở rộng lên $es=3$ (\texttt{Posit(32,3)}), hằng số thang
        $\mathrm{useed}=2^{2^3}=256$ nâng $\mathrm{maxpos}$ lên $2^{240}\approx
        1{,}7668\times 10^{72}$ và hạ $\mathrm{minpos}$ xuống $2^{-240}\approx
        5{,}6598\times 10^{-73}$, mở rộng dải biểu diễn vượt xa FP32 ở cả hai
        phía rất lớn và rất nhỏ.
  \item \textbf{Bảo toàn độ phân giải tích lũy ở nhánh cộng (26 bit fraction quanh $1$):}
        Nhờ đặc tính độ chính xác thuôn dần (\textit{tapered precision}), trong
        miền $1 \le |x| < 2$ ($k=0$, chuỗi regime dài 2 bit), định dạng
        \texttt{Posit(32,3)} cung cấp đến $32 - 1 - 2 - 3 = 26$ bit cho phần
        fraction ($27$ bit significand kể cả bit ẩn), tức bước nhảy lượng tử
        $\Delta x = 2^{-26}$ ở cấp lưu trữ và đường cộng $C$. Đường giải mã
        toán hạng $C$ của \texttt{MAC-F} giữ trọn vẹn $26$ bit này để giảm
        nhiễu làm tròn khi cộng dồn so với tích lũy 16 bit.
  \item \textbf{Đơn giản hóa và nhất quán quy tắc ngoại lệ phần cứng:} Định dạng
        \texttt{Posit(32,3)} chỉ có duy nhất một mã \texttt{Zero} (\texttt{0x00000000})
        và một mã \texttt{NaR} (\texttt{0x80000000}), không có $\pm 0$ và không
        tồn tại số dưới chuẩn (mọi giá trị hữu hạn khác $0$ đều có bit ẩn bằng $1$).
        Các phép toán vượt biên hữu hạn được kẹp bão hòa về $\pm\mathrm{maxpos}$
        hoặc $\pm\mathrm{minpos}$ thay vì tràn thành $\infty$, đồng thời hai số
        Posit có thể so sánh thứ tự trực tiếp như hai số nguyên bù hai 32 bit.
  \item \textbf{Phù hợp với kiến trúc nhân lặp xấp xỉ và tích lũy độ chính xác hỗn hợp:}
        Thay vì triển khai bộ nhân đầy đủ $27\times 27$ bit tiêu tốn nhiều khối DSP,
        bộ nhân lặp \texttt{FULL\_TERM26/Q24} cắt phần fraction đầu vào còn
        $12$ bit ($13$ bit significand kể cả bit ẩn) để đạt chi phí $0\text{ DSP}$,
        trong khi vẫn kế thừa nguyên vẹn băng thông dải động $2^{\pm 240}$ và
        toàn bộ $26$ bit fraction của nhánh cộng $C$ trong lõi hợp nhất \texttt{MAC-F}.
\end{enumerate}

\subsection*{Phân tích sự phối hợp bất đối xứng: vì sao dùng Posit(32,3) với bộ nhân 12 bit?}
Một câu hỏi kiến trúc tự nhiên là: bộ nhân \texttt{FULL\_TERM26/Q24} giới hạn độ
rộng fraction ở $12$ bit, khiến sai số tương đối đơn phép nhân ở mức
$\mathrm{P95}\approx 0{,}0297\%$ ($\approx 2^{-11{,}7}$, tương đương độ chính xác
phần định trị của \texttt{FP16} 10 bit hoặc \texttt{BFloat16} 7 bit); vậy tại sao
hệ thống sử dụng từ mã 32 bit \texttt{Posit(32,3)} thay vì một định dạng 16 bit?
Lựa chọn bất đối xứng này dựa trên ba cơ sở định lượng cho phép nhân--tích lũy
($S_i = S_{i-1} + A_i B_i$):
\begin{enumerate}
  \item \textbf{Nới rộng ngưỡng mất bit tích lũy so với 16 bit và chuẩn hóa RNE đơn phép:}
        Nếu cả tích và thanh ghi cộng $S_{i-1}$ đều giữ ở định dạng 16 bit
        (\texttt{FP16} chỉ có $10$ bit fraction, \texttt{BFloat16} chỉ có $7$ bit),
        khi $S_{i-1}$ lớn hơn $A_i B_i$ quá $2^{11}$ (hoặc $2^8$) lần, toàn bộ
        tín hiệu mới đã bị triệt tiêu hoàn toàn khỏi phần định trị của thanh ghi
        cộng. Ngược lại, khi lưu trữ \texttt{Posit(32,3)} và đưa đủ \textbf{26 bit
        fraction} cho nhánh $C$, ngưỡng chênh lệch thang bắt đầu mất bit cộng
        được nâng thêm tới $2^{16} = 65.536\times$ so với \texttt{FP16} (từ $\Delta sf \approx 11$
        lên $\Delta sf \approx 27$). Bên trong từng phép MAC đơn lẻ, khi $\Delta sf \ge 27$,
        số hạng nhỏ ($< 2^{-27} = 0{,}5\,\mathrm{ULP}$ của số lớn) rơi xuống dưới
        nửa bước lượng tử nên cửa sổ căn chỉnh 31 bit kèm bit \textit{sticky/jam}
        vẫn cho kết quả làm tròn RNE chính xác tuyệt đối. Riêng đối với bài toán
        \emph{chống nuốt số qua chuỗi tích lũy rất dài} nơi tổng chênh vượt biên
        $2^{26}$ lần (như kịch bản \texttt{Cancel}), thiết kế sử dụng thanh ghi
        972 bit của \texttt{MAC-Quire} (Chương~\ref{ch:measurements}).
  \item \textbf{Dải động siêu rộng $2^{\pm 240}$ ($\approx 10^{\pm 72}$) với chi phí 0 DSP:}
        Trong bộ nhân lặp dịch--cộng, độ phức tạp phần cứng chủ yếu tỉ lệ theo số
        bit phần định trị ($12$ bit $\Rightarrow$ tối đa 12 bước lặp, hoàn toàn dùng
        LUT với $0\text{ DSP}$); phần xử lý thang mũ ($sf = sf_A + sf_B$) chỉ là
        một phép cộng số nguyên có dấu 11 bit rất rẻ. Nhờ đó, \texttt{Posit(32,3)}
        mở rộng dải động vượt xa \texttt{IEEE-754 FP32} ($10^{\pm 38}$), \texttt{BFloat16}
        ($10^{\pm 38}$) và \texttt{FP16} ($\pm 65.504$) mà không làm phình bộ nhân.
  \item \textbf{Miền ứng dụng mục tiêu:} Mô hình tích lũy độ chính xác hỗn hợp
        (nhân xấp xỉ $\approx 12$ bit định trị, cộng trên lưới 26--31 bit dải
        động siêu rộng) nhắm trực tiếp tới suy luận mạng nơ-ron sâu trên FPGA biên
        (chấp nhận nhiễu trọng số 8--12 bit), bộ lọc số FIR/IIR băng thông vừa và
        tích vô hướng tín hiệu cảm biến có biên độ trải rộng nhiều bậc độ lớn.
\end{enumerate}

\begin{engineeringtable}{So sánh định lượng giữa chuẩn IEEE-754 binary32, Posit(32,2) và Posit(32,3).}{tab:ieee-vs-posit}
\begin{tabularx}{\linewidth}{@{}>{\raggedright\arraybackslash}p{42mm} Y Y Y@{}}
\toprule
Đặc tính & IEEE-754 binary32 & Posit(32,2) & Posit(32,3) (đã chọn) \\
\midrule
Độ rộng từ mã & 32 bit & 32 bit & 32 bit \\
Số bit exponent ($e$) & Cố định 8 bit ($bias=127$) & Tối đa 2 bit ($\mathrm{useed}=16$) & Tối đa 3 bit ($\mathrm{useed}=256$) \\
Fraction quanh $|x|\approx 1$ & Cố định 23 bit & 27 bit ($k=0$) & 26 bit ($k=0$) \\
Bước nhảy $\Delta x$ trong $[1,2)$ & $2^{-23}$ & $2^{-27}$ (mịn $16\times$) & $2^{-26}$ (mịn $8\times$) \\
$\mathrm{minpos}$ hữu hạn & $2^{-149}$ (subnormal) / $2^{-126}$ (normal) & $2^{-120}\approx 7{,}52\times 10^{-37}$ & $2^{-240}\approx 5{,}66\times 10^{-73}$ \\
$\mathrm{maxpos}$ hữu hạn & $(2-2^{-23})2^{127}\approx 3{,}40\times 10^{38}$ & $2^{120}\approx 1{,}33\times 10^{36}$ & $2^{240}\approx 1{,}77\times 10^{72}$ \\
Số dưới chuẩn (subnormal) & Có (bit ẩn $=0$) & Không (bit ẩn luôn $=1$) & Không (bit ẩn luôn $=1$) \\
Mã đặc biệt & $\pm 0, \pm\infty$, nhiều \texttt{NaN} & 1 \texttt{Zero}, 1 \texttt{NaR} & 1 \texttt{Zero}, 1 \texttt{NaR} \\
\bottomrule
\end{tabularx}
\end{engineeringtable}

\section{Định dạng Posit(32,3) và quy tắc giải mã}
\begin{figure}[htbp]
  \centering
  \begin{tikzpicture}[x=1cm,y=1cm]
    \fill[blue!9] (0,0) rectangle (1,0.85);
    \fill[blue!4] (1,0) rectangle (4,0.85);
    \fill[orange!14] (4,0) rectangle (6.7,0.85);
    \fill[blue!4] (6.7,0) rectangle (11.5,0.85);
    \draw (0,0) rectangle (11.5,0.85);
    \foreach \x in {1,4,6.7} \draw (\x,0)--(\x,0.85);
    \node at (0.5,0.425) {$s$};
    \node at (2.5,0.425) {Regime};
    \node at (5.35,0.425) {Exponent};
    \node at (9.1,0.425) {Fraction};
    \node[font=\scriptsize] at (0.5,-.22) {1 bit};
    \node[font=\scriptsize] at (2.5,-.22) {Độ dài thay đổi};
    \node[font=\scriptsize] at (5.35,-.22) {Tối đa 3 bit};
    \node[font=\scriptsize] at (9.1,-.22) {Các bit còn lại};
  \end{tikzpicture}
  \caption{Các trường của độ lớn hữu hạn Posit(32,3) sau khi xử lý dấu.}
  \label{fig:posit-fields}
\end{figure}

Một từ mã \texttt{Posit(32,3)} hữu hạn khác 0 bao gồm 4 thành phần theo thứ tự từ MSB đến LSB: dấu ($s$), \textit{regime} ($r$), \textit{exponent} ($e$) và \textit{fraction} ($f$) \cite{gustafson2017posit,positstandard2022}. Hai mã đặc biệt được nhận diện ngay từ từ mã thô trước khi giải mã các trường: \texttt{0x00000000} biểu diễn chính xác số $0$ (\texttt{Zero}) và \texttt{0x80000000} biểu diễn giá trị ngoại lệ (\texttt{NaR} -- \textit{Not a Real}). Với mọi mã âm hữu hạn ($s=1$, khác \texttt{NaR}), quy tắc chuẩn Posit yêu cầu lấy bù hai toàn bộ 32 bit trước khi tách trường:
\[
P_{\mathrm{mag}} = \begin{cases} P, & s = 0, \\ (-P) \bmod 2^{32} = \overline{P} + 1 \pmod{2^{32}}, & s = 1. \end{cases}
\]

Trên 31 bit độ lớn $P_{\mathrm{mag}}[30:0]$, trường \textit{regime} là một chuỗi $m \ge 1$ bit giống nhau liên tiếp kết thúc bằng một bit đối lập (\textit{terminator}) hoặc chạm tới hết từ mã. Nếu chuỗi mở đầu gồm $m$ bit~1 (kết thúc bằng 0) thì $k = m - 1 \ge 0$; nếu chuỗi mở đầu gồm $m$ bit~0 (kết thúc bằng 1) thì $k = -m < 0$. Đóng góp thang của regime là $\mathrm{useed}^k = 2^{2^{es}\cdot k} = 2^{8k}$ (Bảng~\ref{tab:regime-rules}).

\begin{engineeringtable}{Quy tắc giải mã trường regime trong Posit(32,3) ($\mathrm{useed}=256$).}{tab:regime-rules}
\begin{tabularx}{\linewidth}{@{}lccc Y@{}}
\toprule
Chuỗi bit regime & Số bit lặp ($m$) & Bit kết thúc & Giá trị $k$ & Hệ số thang $\mathrm{useed}^k$ \\
\midrule
\texttt{0001} & 3 & 1 & $-3$ & $256^{-3} = 2^{-24}$ \\
\texttt{001} & 2 & 1 & $-2$ & $256^{-2} = 2^{-16}$ \\
\texttt{01} & 1 & 1 & $-1$ & $256^{-1} = 2^{-8}$ \\
\texttt{10} & 1 & 0 & $0$ & $256^{0} = 1$ \\
\texttt{110} & 2 & 0 & $1$ & $256^{1} = 2^{8}$ \\
\texttt{1110} & 3 & 0 & $2$ & $256^{2} = 2^{16}$ \\
\bottomrule
\end{tabularx}
\end{engineeringtable}

Sau bit kết thúc của regime, tối đa $es=3$ bit tiếp theo tạo thành số mũ không dấu $e \in [0, 7]$ (nếu số bit còn lại ít hơn 3, các bit thấp còn thiếu được đệm bằng 0). Các bit còn lại (nếu có) tạo thành phần phân số $f \in [0, 1)$, kết hợp với bit ẩn luôn bằng 1 tạo thành significand chuẩn hóa $1 + f \in [1, 2)$. Đặt số mũ tổng hợp (\textit{scale factor}) là $sf = 8k + e$, giá trị số thực của một mã \texttt{Posit(32,3)} hữu hạn khác 0 được xác định thống nhất bởi:
\[
x = (-1)^s \cdot \mathrm{useed}^k \cdot 2^e \cdot (1 + f) = (-1)^s \cdot 2^{8k+e} \cdot (1 + f) = (-1)^s \cdot 2^{sf} \cdot (1 + f).
\]

\subsection{Dấu, mã độ lớn và miền giá trị}

Đối với số âm hữu hạn, phần cứng thực hiện lấy bù hai 32 bit trước khi đếm chuỗi
regime; chẳng hạn từ mã \texttt{0xbc000000} được đổi thành độ lớn \texttt{0x44000000},
giải mã phần trường như $+2$ rồi khôi phục dấu âm ở đầu ra. Nếu đếm trực tiếp trên
mã âm thô, chuỗi bit lặp sẽ phản ánh dạng mã hóa bù hai chứ không còn mang ý nghĩa
của trường regime. Hai mã \texttt{Zero} và \texttt{NaR} được kiểm tra trước bước
tách trường để tránh nhầm lẫn cực trị hữu hạn với trạng thái ngoại lệ.

\begin{engineeringtable}{Các mã tiêu biểu của Posit(32,3) dùng để đối chiếu giải mã và các cận biên.}{tab:posit-code-examples}
\begin{tabularx}{\linewidth}{@{}llllY@{}}
\toprule
Mã hex & Dấu & $k$ & $e$ & Giá trị hoặc vai trò \\
\midrule
\texttt{00000000} & 0 & -- & -- & \texttt{Zero} ($0$) \\
\texttt{80000000} & 1 & -- & -- & \texttt{NaR} (\textit{Not a Real}) \\
\texttt{40000000} & 0 & $0$ & $0$ & $1$ \\
\texttt{44000000} & 0 & $0$ & $1$ & $2$ \\
\texttt{bc000000} & 1 & $0$ & $1$ & $-2$ sau khi lấy bù hai thành \texttt{44000000} \\
\texttt{00000001} & 0 & $-30$ & $0$ & $2^{-240}$ ($\mathrm{minpos}$) \\
\texttt{7fffffff} & 0 & $30$ & $0$ & $2^{240}$ ($\mathrm{maxpos}$) \\
\bottomrule
\end{tabularx}
\end{engineeringtable}

Khi số bit còn lại sau regime không đủ 3 bit exponent, các vị trí thấp còn thiếu
được đệm 0 theo quy tắc giải mã Posit. Không gian dành cho fraction thu hẹp dần khi
chuỗi regime kéo dài về phía gần $\mathrm{minpos}$ hoặc $\mathrm{maxpos}$; ngược
lại quanh giá trị $1$, regime chỉ dài 2 bit và để lại tới 26 bit fraction. Do bước
lượng tử biến thiên theo từng miền $k$, không thể xem hiệu giữa hai mã Posit liên
tiếp là một đơn vị ULP có độ lớn không đổi trên toàn dải số.

\subsection{Ví dụ giải mã Posit(32,3) dương và âm}

\paragraph{Ví dụ 1: Giải mã một số Posit dương.}
Xét từ mã đầu vào $P_+ = \texttt{0x1B400000}$, có biểu diễn nhị phân 32 bit là:
\[
P_+ = \texttt{00011011\ 01000000\ 00000000\ 00000000}_2.
\]
Vì bit MSB bằng $0$ và từ mã khác \texttt{0x00000000}, đây là một số dương hữu
hạn khác $0$. Các trường bit tính từ trái sang phải được tách thành 4 phần (Bảng~\ref{tab:decode-ex-pos}).

\begin{engineeringtable}{Phân tách trường bit cho ví dụ giải mã số Posit dương \texttt{0x1B400000}.}{tab:decode-ex-pos}
\begin{tabularx}{\linewidth}{@{}>{\raggedright\arraybackslash}p{26mm} l r Y@{}}
\toprule
Trường & Chuỗi bit & Số bit & Ý nghĩa giải mã \\
\midrule
Sign ($s$) & \texttt{0} & 1 & Số dương ($s=0$). \\
Regime & \texttt{001} & 3 & Hai bit \texttt{0} liên tiếp rồi kết thúc bằng \texttt{1}: $m=2 \Rightarrow k=-m=-2$. \\
Exponent ($e$) & \texttt{101} & 3 & Giá trị không dấu $e = 101_2 = 5$. \\
Fraction ($f$) & \texttt{101000...00} & 25 & $f = 1\cdot 2^{-1} + 0\cdot 2^{-2} + 1\cdot 2^{-3} = \frac{5}{8}$. \\
\bottomrule
\end{tabularx}
\end{engineeringtable}

Từ các trường vừa tách, hệ số tỷ lệ của regime là $\mathrm{useed}^k = 256^{-2} = 2^{-16}$,
phần trị có bit ẩn bằng 1 là $M = 1 + f = 1 + \frac{5}{8} = \frac{13}{8} = 1{,}625$,
và số mũ tổng hợp bằng $sf = 8k + e = 8(-2) + 5 = -11$. Giá trị thực tương ứng của
từ mã là:
\[
x_+ = (-1)^0 \cdot 2^{-11} \cdot \frac{13}{8} = \frac{13}{16\,384} = 0{,}00079345703125.
\]

\paragraph{Ví dụ 2: Giải mã một số Posit âm.}
Xét từ mã đối của ví dụ trên là $P_- = \texttt{0xE4C00000}$:
\[
P_- = \texttt{11100100\ 11000000\ 00000000\ 00000000}_2.
\]
Bit dấu $s=1$ và mã khác \texttt{0x80000000} nên đây là một số âm hữu hạn. Phần
cứng lưu dấu $s=1$ rồi thực hiện bù hai toàn bộ 32 bit để đưa về mã độ lớn:
\[
P_{\mathrm{mag}} = \overline{\texttt{0xE4C00000}} + 1 = \texttt{0x1B3FFFFF} + 1 = \texttt{0x1B400000}.
\]
Tách trường trên $P_{\mathrm{mag}}$ thu được cùng bộ tham số $k=-2$, $e=5$,
$M=\frac{13}{8}$ và $sf=-11$ như Ví dụ 1. Ghép với dấu $s=1$ ban đầu cho kết quả:
\[
x_- = (-1)^1 \cdot 2^{-11} \cdot \frac{13}{8} = -\frac{13}{16\,384} = -0{,}00079345703125.
\]
Hai ví dụ cho thấy việc tách trường trên số âm bắt buộc phải thực hiện sau bước
lấy bù hai; đồng thời cần phân biệt rõ bit dấu $s$ của toán hạng với dấu âm/dương
của các đại lượng nội bộ $k$ và $sf$.

\subsection{Lưới số nguyên bảo toàn trọng số bit}

Với phần fraction được căn đủ 26 bit ở định dạng \texttt{Posit(32,3)}, hệ số nguyên
không dấu mang bit ẩn là $J = 2^{26} + F$, nên mọi số hữu hạn khác 0 biểu diễn
chính xác dưới dạng $x = (-1)^s J 2^{sf - 26}$. Khi nhân hai toán hạng bất kỳ,
$sf_A + sf_B \ge -480$ và tích $J_A J_B$ có $52$ vị trí fraction; do đó toàn bộ
trọng số thấp nhất của tích nằm trên lưới số nguyên chung $2^{-532}$. Biểu diễn
bằng số nguyên lớn trên lưới $2^{-532}$ được mô hình tham chiếu (\textit{golden
model}) sử dụng để thực hiện phép cộng và đánh giá sai số độc lập hoàn toàn với
giới hạn độ rộng của dấu phẩy động \texttt{binary64}.

Bên trong phần cứng, dữ liệu không cần duy trì lưới rộng ở mọi tầng: khối nhân
lặp làm việc trên lưới 24 bit sau dấu chấm (\texttt{Q24}), bộ cộng của \texttt{MAC-F}
làm việc trên cửa sổ 31 bit (\texttt{Q29}) kèm bit phần dư căn chỉnh, còn
\texttt{MAC-Quire} duy trì trạng thái 972 bit trên lưới $2^{-480}$ để tích lũy
liên tiếp nhiều tích xấp xỉ trước lệnh \texttt{READ}.


\section{Nguyên lý nhân lặp}
Kế thừa phương pháp nhân Posit lặp của Norris và Kim \cite{norris2021iterativeposit},
bộ nhân sử dụng ba khối chức năng: khối chọn toán hạng \texttt{OPS}, khối tính lượng
dịch \texttt{SAC} tìm vị trí bit 1 tiếp theo, và khối cộng dịch \texttt{SBM} để
cộng dồn từng số hạng. Ở ngõ vào, bộ nhân giữ tối đa $12$ bit cao của fraction
từng toán hạng (đệm 0 đối với fraction ngắn hơn $12$ bit). Nếu $S_i \in \{1,\dots,12\}$
là vị trí tính từ dấu chấm nhị phân của các bit 1 trong $f_X$, phép nhân hai
significand chuẩn hóa được khai triển thành tổng các bản dịch phải của $1+f_Y$:
\[
    (1+f_X)(1+f_Y) = (1+f_Y) + \sum_{i} (1+f_Y) 2^{-S_i}.
\]

Trong RTL của đề tài, bộ nhân xử lý tối đa \texttt{cfg\_n} $\in [0, 12]$ số hạng
fraction theo thứ tự từ trọng số lớn đến trọng số nhỏ và tự động dừng sớm ngay khi
mặt nạ các bit 1 còn lại bằng 0. Ví dụ với $f_X = 0.01001_2$, chỉ cần tối đa hai
vòng lặp tương ứng với $S_1 = 2$ và $S_2 = 5$ miễn là $\texttt{cfg\_n} \ge 2$. Cả hai
nguồn xấp xỉ gồm cắt fraction đầu vào về 12 bit và dừng vòng lặp sớm theo \texttt{cfg\_n}
đều được đánh dấu bằng hai cờ trạng thái tách biệt (\texttt{input\_cut} và
\texttt{approx\_cut}).

\section{Hợp đồng phép nhân--cộng}
Ký hiệu $\widehat P$ là tích nội bộ trên lưới \texttt{FULL\_TERM26/Q24} sinh ra ở
đầu ra bộ nhân lặp, và $R(\cdot)$ là toán tử chuẩn hóa, làm tròn RNE và đóng gói
về từ mã \texttt{Posit(32,3)} theo quy ước làm tròn chuẩn của đề tài. Đối với các toán
hạng đầu vào hữu hạn thông thường, hai biến thể kiến trúc thực hiện hai ánh xạ số
học sau:
\[
  D_{\mathrm{NF}} = R\bigl(R(\widehat P) + C\bigr), \qquad
  D_{\mathrm{F}}  = R(\widehat P + C).
\]

Bên cạnh hai lõi trực tiếp nhận toán hạng $C$ từ ngoài, kiến trúc đối chứng
\texttt{MAC-Quire} duy trì bộ tích lũy có dấu 972 bit trên lưới $2^{-480}$, hỗ
trợ tích lũy tối đa 1.024 phép MAC liên tiếp giữa hai lệnh \texttt{CLEAR}. Lệnh
\texttt{MAC} cộng trực tiếp tích $\widehat P$ lên lưới $2^{-480}$ của thanh ghi
$Q$, còn lệnh \texttt{READ} thực hiện $R(Q)$ để trả kết quả \texttt{Posit(32,3)}
mà vẫn bảo toàn trạng thái $Q$.

Cả \texttt{MAC-NF} và \texttt{MAC-F} cùng nhận bộ dữ liệu \texttt{a/b/c[31:0]},
cấu hình \texttt{cfg\_n[3:0]} và \texttt{cfg\_ops[1:0]} tại cạnh bắt tay đầu vào,
và xuất kết quả \texttt{d[31:0]} cùng thanh ghi 5 cờ \texttt{flags[4:0]} theo thứ
tự \{\texttt{nar}, \texttt{sat\_max}, \texttt{sat\_min}, \texttt{inexact},
\texttt{approx\_cut}\}. Trong chế độ mặc định \texttt{OPS=1/n=12}, \texttt{OPS}
chọn toán hạng có ít bit 1 hơn trong 12 bit fraction làm $X$ (khi số bit 1 bằng
nhau thì giữ $A$ làm $X$); ngược lại \texttt{OPS=0} luôn cố định $A$ làm $X$.

\section{Độ rộng tích, sai số và chuẩn hóa phép nhân}

Đặt $Y = 4096 + F_Y$ là significand nguyên 13 bit (gồm bit ẩn và 12 bit fraction
cao) của toán hạng được dịch, và $M_0 = Y 2^{12}$ là giá trị nền khởi tạo trên
lưới \texttt{Q24}. Vì $M_0$ có $12$ bit thấp bằng $0$, mọi lượng dịch phải
$S_i \in \{1, \dots, 12\}$ đều bảo toàn toàn bộ trọng số khác $0$ của từng số
hạng tích riêng phần. Khi xử lý trọn vẹn mọi bit $1$ trong 12 bit fraction của $X$,
tích nguyên bằng $(4096+F_X)(4096+F_Y)$ với cận trên $8191^2 = 67\,092\,481 < 2^{26}$.
Do đó 26 bit không dấu của thanh ghi tích lũy \texttt{Q24} (2 bit nguyên và 24 bit
fraction) đủ chứa mọi tích có thể xảy ra mà không tràn ngoài miền hợp lệ; bit 25
bằng $1$ chính là tín hiệu phát hiện tràn chuẩn hóa ($M/2^{24} \ge 2$).

Gọi $a = 1+f_A$ và $b = 1+f_B$ là hai significand đầy đủ ban đầu, còn
$a' = a - \delta_a$ và $b' = b - \delta_b$ là hai significand sau khi giữ 12 bit
fraction cao ($0 \le \delta_a, \delta_b < 2^{-12}$). Độ lệch do bước cắt đầu vào
được phân tách đại số theo biểu thức:
\[
ab - a'b' = a'\delta_b + b'\delta_a + \delta_a\delta_b.
\]
Từ tính chất không âm của các phần dư $\delta_a, \delta_b$ và các số hạng dịch
tiếp theo, với cùng một cặp đầu vào $(A, B)$, ký hiệu độ lớn của tích đầy đủ,
tích 12 bit fraction và tích RTL ở giới hạn $n$ vòng lần lượt là $|P|$, $|P_{12}|$
và $|\widehat P|$, ta luôn có tính đơn điệu theo độ lớn tại đầu ra bộ nhân:
\[
|\widehat P| \le |P_{12}| \le |P|.
\]
Hệ quả quan trọng là bộ nhân luôn mang \textbf{thiên lệch cắt một phía về không}
(\textit{one-sided truncation bias}): phép kiểm tra toàn diện trên toàn bộ
$125.668$ cặp $(A,B)$ hữu hạn khác $0$ của tập mẫu thực nghiệm ở cấu hình
\texttt{OPS1/n12} ghi nhận bất đẳng thức $|\widehat P| \le |P|$ đúng tuyệt đối
trên \textbf{100\% mẫu} ($117.486$ cặp nhỏ hơn chặt và $8.182$ cặp bằng nhau
chính xác khi 14 bit fraction thấp đều bằng $0$). Khi cộng với toán hạng trái dấu
$C$ và đi qua bước lượng tử hóa RNE ở cuối MAC, thiên lệch âm có hệ thống này đôi
khi được bù trừ ngẫu nhiên bởi chiều làm tròn lên của bước đóng gói trung gian
$R(\widehat P)$ trong \texttt{MAC-NF} (giải thích hiện tượng một số phân vị
tương đối đơn bước tại Chương~\ref{ch:measurements}).

\begin{engineeringtable}{Phân tách các nguồn sai số và mốc đối chiếu độc lập.}{tab:error-mechanisms}
\begin{tabularx}{\linewidth}{@{}>{\raggedright\arraybackslash}p{32mm}Y Y@{}}
\toprule
Nguồn sai số & Cơ chế và cờ trạng thái trong RTL & Mốc đối chiếu độc lập \\
\midrule
Cắt fraction đầu vào & Loại bỏ bit thấp ngoài 12 bit ở A hoặc B; chốt cờ \texttt{input\_cut}. & So tích 12 bit $P_{12}$ với tích chính xác $P$ đủ 26 bit fraction. \\
Giới hạn vòng lặp & Mặt nạ $X$ vẫn còn bit 1 khi dừng tại $t=n$; xuất cờ \texttt{approx\_cut}. & So tích RTL $\widehat P$ với $P_{12}$ tại \texttt{core\_done}. \\
Lượng tử hóa khi đóng gói & Giá trị giải mã từ từ mã 32 bit lệch so với giá trị trước làm tròn RNE. & So $\operatorname{decode}(D)$ với $R(\widehat P + C)$ (hoặc $R(R(\widehat P)+C)$). \\
\bottomrule
\end{tabularx}
\end{engineeringtable}

\subsection{Ví dụ minh họa phép nhân Posit(32,3) từng bước}

Xét phép nhân hai số dương $A = 3{,}5 = +2^1 \times 1{,}75$ (mã \texttt{0x47000000})
và $B = 0{,}75 = +2^{-1} \times 1{,}5$ (mã \texttt{0x3E000000}). Trình tự 4 bước
thực hiện phép nhân Posit được minh họa chi tiết như sau:

\begin{enumerate}
  \item \textbf{Bước 1 (Giải mã đầu vào):}
        $A$ cho $s_A = 0$, $sf_A = 1$, $M_A = 1.1100_2$;
        $B$ cho $s_B = 0$, $sf_B = -1$, $M_B = 1.1000_2$.
  \item \textbf{Bước 2 (Tính dấu và scale thô):}
        $s_P = s_A \oplus s_B = 0$ (dương) và $sf_{\mathrm{raw}} = sf_A + sf_B = 1 + (-1) = 0$.
  \item \textbf{Bước 3 (Nhân phần trị theo cột nhị phân):}
        Thực hiện nhân hai significand trên lưới bốn vị trí fraction:
        \[
        \begin{array}{r}
          \phantom{\times}01.1100\\
          \underline{\times\phantom{0}01.1000}\\
          \phantom{+}00.0000\\
          \phantom{+}00.1110\\
          \underline{+\phantom{0}01.1100}\\
          \phantom{+}10.1010
        \end{array}
        \]
        Thu được tích thô $M_{\mathrm{raw}} = 10.1010_2 = 2{,}625_{10}$.
  \item \textbf{Bước 4 (Chuẩn hóa sau nhân và đóng gói):}
        Vì $M_{\mathrm{raw}} = 10.1010_2 \ge 2$ (bit nguyên tràn sang vị trí $2^1$),
        dịch phải 1 bit để đưa bit 1 cao nhất về vị trí ẩn: $M_P = 1.0101_2$ (bỏ bit
        ẩn còn fraction \texttt{0101}), đồng thời tăng scale: $sf_P = sf_{\mathrm{raw}} + 1 = 1$.
        Khi đóng gói, $sf_P = 1 \Rightarrow k = \lfloor 1/8 \rfloor = 0$ (regime
        \texttt{10}) và $e = 1 = \texttt{001}_2$. Ghép trường thu được chuỗi bit
        \texttt{0\_10\_001\_01010...} $= \texttt{0x45400000}$, biểu diễn đúng giá
        trị $2^1 \times 1{,}3125 = 2{,}625_{10}$.
\end{enumerate}

\subsection{Ví dụ khai triển tích lặp trên lưới FULL\_TERM26/Q24}

Để minh họa trực tiếp cơ chế dịch--cộng theo bit trong RTL của đề tài, xét hai
significand $M_X = 1.1010_2 = 1{,}625$ và $M_Y = 1.1000_2 = 1{,}5$. Fraction của
$X$ chỉ có hai bit $1$ tại các vị trí $S_1 = 1$ và $S_2 = 3$ sau dấu chấm nhị phân.
Với quy ước 12 bit fraction, $Y = 4096 + 2048 = 6144$ và giá trị khởi tạo lưới
Q24 là $M_0 = 6144 \times 4096 = 25\,165\,824$. Khai triển tuần tự theo vòng lặp:
\[
1.1010_2 \times 1.1000_2 = 1.1000_2 + 0.1100_2 + 0.0011_2 = 10.0111_2 = 2{,}4375_{10}.
\]
Ở cấp số nguyên Q24 của phần cứng:
\begin{itemize}
  \item Khởi tạo (\texttt{first}, đóng góp bit ẩn): $M_0 = 25\,165\,824$ (tương ứng $M_0/2^{24} = 1{,}5$ nếu $n=0$).
  \item Vòng 1 ($S_1 = 1$): $T_1 = 6144 \times 2^{12-1} = 12\,582\,912 \Rightarrow M_1 = M_0 + T_1 = 37\,748\,736$ (tương ứng $2{,}25$ nếu $n=1$).
  \item Vòng 2 ($S_2 = 3$): $T_2 = 6144 \times 2^{12-3} = 3\,145\,728 \Rightarrow M_2 = M_1 + T_2 = 40\,894\,464$, suy ra $\frac{M_2}{2^{24}} = 2{,}4375$. Mặt nạ của $X$ cạn về $0$ nên lõi tự động dừng sớm ở $t=2$ vòng mà không chạy các vòng rỗng khi cấu hình $n=12$.
\end{itemize}

\paragraph{Đối chiếu trên cặp mã 32 bit FT-05.}
Với bộ dữ liệu kiểm thử RTL \texttt{FT-05}: $A = \texttt{0x1c900000}$ (regime
\texttt{001}, $k_A=-2$, $e_A=6 \Rightarrow sf_A=-10$, $F_A = \texttt{0x480}$) và
$B = \texttt{0x3c820000}$ (regime \texttt{01}, $k_B=-1$, $e_B=7 \Rightarrow sf_B=-1$,
$F_B = \texttt{0x208}$). Cả hai mã đều có 2 bit 1 trong 12 bit fraction cao và
14 bit thấp bằng 0, nên \texttt{OPS1} giữ $A$ làm $X$ và $B$ làm $Y$, cho $b = sf_A + sf_B = -11$.
Quá trình cộng trên lưới Q24 diễn ra như Bảng~\ref{tab:ft05-worked}. Giá trị trước
Packer là $24\,224\,768 \times 2^{-35}$; vì bit 25 bằng 0 nên adapter giữ nguyên
$sf = -11$. Khi $C=0$, cả mô hình tham chiếu và RTL cùng xuất mã \texttt{0x1ae34800}
(khi đặt $n=0$, bộ nhân chỉ phát đóng góp cơ số $M_0 = 18\,907\,136$ và bật cờ
\texttt{approx\_cut=1} để cho ra mã \texttt{0x1a410000}).

\begin{engineeringtable}{Phép tính FULL\_TERM26/Q24 theo hệ số nguyên cho cặp FT-05.}{tab:ft05-worked}
\begin{tabularx}{\linewidth}{@{}>{\raggedright\arraybackslash}p{34mm}r Y@{}}
\toprule
Đại lượng & Giá trị nguyên & Diễn giải nguồn gốc \\
\midrule
$Y$ & 4.616 & Bit ẩn $4.096$ cộng fraction $B$ ($F_B = 520$). \\
$M_0$ & 18.907.136 & Cơ số $Y \ll 12$ tại xung \texttt{first}. \\
Term 1 ($S_1=2$) & 4.726.784 & Bit 1 đầu tiên của $X$, $M_0 \gg 2$. \\
Term 2 ($S_2=5$) & 590.848 & Bit 1 thứ hai của $X$, $M_0 \gg 5$. \\
$M$ tại \texttt{last} & 24.224.768 & Tổng cơ số và hai term, $t=2$, $b=-11$. \\
\bottomrule
\end{tabularx}
\end{engineeringtable}

\section{Nguyên lý cộng, căn chỉnh, jam và chuẩn hóa triệt tiêu}

Khi cộng hoặc trừ hai số hạng có dấu, scale và significand đã chuẩn hóa, bước
đầu tiên là so sánh độ lớn (so $sf$ trước, nếu bằng nhau thì so significand) để
đặt số hạng lớn hơn làm toán hạng nền $X$ và số hạng còn lại làm $Y$. Khi hai
toán hạng cùng dấu hiệu dụng, thực hiện phép cộng $M_X + (M_Y \gg \Delta sf)$;
khi khác dấu hiệu dụng, thực hiện phép trừ không âm $M_X - (M_Y \gg \Delta sf)$
và giữ dấu của $X$.

\paragraph{Ví dụ minh họa phép cộng phát sinh carry và phép trừ triệt tiêu.}
Xét cặp toán hạng $A = 3{,}5$ (\texttt{0x47000000}, $sf_A = 1$, $M_A = 1.1100_2$)
và $B = 0{,}75$ (\texttt{0x3E000000}, $sf_B = -1$, $M_B = 1.1000_2$):
\begin{enumerate}
  \item \textbf{Phép cộng có nhớ ($A + B$):} Vì $sf_A > sf_B$, chọn $A$ làm nền
        với độ lệch $\Delta sf = 1 - (-1) = 2$. Dịch phải $M_B$ hai vị trí thu
        được $0.0110_2$ trên cùng thang $2^1$:
        \[
        3{,}5 + 0{,}75 = 2^1\bigl(1.1100_2 + 0.0110_2\bigr),
        \qquad
        \begin{array}{r}
          \phantom{+}01.1100\\
          \underline{+\phantom{0}00.0110}\\
          \phantom{+}10.0010
        \end{array}
        \]
        Bit nhớ tràn vào vị trí nguyên thứ hai tạo tổng thô $10.0010_2$. Bước chuẩn
        hóa sau carry dịch phải 1 bit và tăng scale thêm $1$:
        $10.0010_2 \times 2^1 = 1.0001_2 \times 2^2 = 4{,}25$. Mã hóa $k=0$, $e=2$
        (\texttt{010}), fraction \texttt{0001} cho ra mã kết quả \texttt{0x48400000}.
  \item \textbf{Phép trừ gây triệt tiêu ($A - 3{,}25$):} Thay số hạng thứ hai
        bằng $-3{,}25$ ($s=1$, $sf=1$, $M = 1.1010_2$). Do cùng $sf=1$ và khác
        dấu, phép cộng chuyển thành trừ hai độ lớn:
        \[
        \begin{array}{r}
          \phantom{-}01.1100\\
          \underline{-\phantom{0}01.1010}\\
          \phantom{-}00.0010
        \end{array}
        \qquad\Longrightarrow\qquad
        0.0010_2 \times 2^1 = 1.0000_2 \times 2^{-2} = 0{,}25.
        \]
        Hiệu độ lớn mất ba bit cao, buộc mạch chuẩn hóa phải dịch trái $3$ vị
        trí và giảm scale từ $1$ xuống $-2$, đóng gói thành \texttt{0x38000000}.
        Trong cửa sổ cộng 31 bit có vị trí bit ẩn $h=29$, ký hiệu $j$ là chỉ số
        bit 1 cao nhất của tổng/hiệu và $\ell = 30 - j$ là số bit 0 dẫn đầu đếm
        từ bit 30, công thức cập nhật scale hợp nhất cho cả trường hợp có nhớ
        ($\ell = 0$), bình thường ($\ell = 1$) và triệt tiêu ($\ell > 1$) là:
        \[
        sf_{\mathrm{out}} = sf_{\mathrm{base}} + j - h = sf_{\mathrm{base}} + 1 - \ell.
        \]
        Áp dụng cho ví dụ trừ trên với $sf_{\mathrm{base}}=1$ và $\ell=4$ ($j=26$)
        thu được đúng $sf_{\mathrm{out}} = 1 + 1 - 4 = -2$. Nếu hai độ lớn triệt
        tiêu hoàn toàn bằng $0$, cờ \texttt{Zero} được kích hoạt trực tiếp thay
        cho kết quả đếm của LOD.
\end{enumerate}

\paragraph{Phần dư căn chỉnh và thao tác jam.}
Trong cửa sổ cộng độ rộng hữu hạn, các bit của toán hạng nhỏ $Y$ bị dịch ra ngoài
cạnh phải của cửa sổ được gom bằng phép OR (\textit{sticky}). Nếu giữ bit LSB
của cửa sổ phản ánh phần dư khác 0 (gọi là \textit{jam}), mạch biết được độ lớn
thật của số hạng nhỏ lớn hơn phần còn giữ trong cửa sổ. Trong phép trừ, phần dư
này làm giảm nhẹ độ lớn hiệu và ảnh hưởng trực tiếp tới biên làm tròn khi xảy ra
triệt tiêu; vì vậy đường dữ liệu phải truyền thông tin phần dư qua cả tầng số học
và tầng chuẩn hóa đến Packer thay vì chỉ OR mù một cờ mất dữ liệu ở cuối.

\section{Làm tròn RNE, mã hóa đầu ra và ranh giới ba kiến trúc}

Sau chuẩn hóa, từ $sf_{\mathrm{out}}$ tách $k = \lfloor sf_{\mathrm{out}} / 8 \rfloor$
và $e = sf_{\mathrm{out}} - 8k \in [0, 7]$ bằng phép chia lấy phần nguyên xuống
(ví dụ $sf = -11 \Rightarrow k = -2, e = 5$). Chuỗi bit $1+\text{terminator}$ hoặc
$0+\text{terminator}$ của $k$ ghép với exponent và fraction tạo thành vector độ
lớn được giữ lại $31$ bit. Gọi $L$ là bit thấp nhất còn giữ, $G$ (\textit{Guard})
là bit kề ngay sau $L$, $R$ (\textit{Round}) là bit tiếp theo và $S$ (\textit{Sticky})
là phép OR toàn bộ các bit thấp hơn cùng phần dư căn chỉnh, điều kiện tăng mã
độ lớn theo quy tắc RNE (\textit{Round to Nearest, Ties to Even}) là:
\[
\rho = G \land (L \lor R \lor S).
\]

\begin{engineeringtable}{Bảng quyết định làm tròn RNE theo các bit L, G, R và S.}{tab:rne-decision}
\begin{tabularx}{\linewidth}{@{}llllY@{}}
\toprule
$G$ & $R$ & $S$ & $L$ & Hành động làm tròn \\
\midrule
0 & bất kỳ & bất kỳ & bất kỳ & Làm tròn xuống (giữ nguyên độ lớn). \\
1 & 0 & 0 & 0 & Hòa chẵn (\textit{tie-even}): giữ mã có $L=0$. \\
1 & 0 & 0 & 1 & Hòa lẻ (\textit{tie-odd}): tăng 1 đơn vị để chuyển $L$ về $0$. \\
1 & 1 & bất kỳ & bất kỳ & Vượt nửa bước: tăng độ lớn thêm 1 đơn vị. \\
1 & 0 & 1 & bất kỳ & Vượt nửa bước (do có dư ở $S$): tăng độ lớn thêm 1 đơn vị. \\
\bottomrule
\end{tabularx}
\end{engineeringtable}

\paragraph{Ví dụ hợp nhất làm tròn RNE và mã hóa bù hai cho số âm.}
Giả sử phần độ lớn giữ được đang là $m = 6$ và quyết định RNE là tăng $\rho = 1$.
Nếu kết quả dương, đầu ra là $m + \rho = 7$; nếu kết quả mang dấu âm, giá trị cần
xuất là $-(m+\rho) = -7$. Trên từ mã 32 bit, đồng nhất thức đại số
\[
-(m + \rho) = \overline{m} + (1 - \rho) \pmod{2^{32}}
\]
cho thấy thay vì cộng xong rồi mới lấy bù hai bằng hai bộ cộng nối tiếp, phần
cứng chỉ cần đảo bit $m$ theo dấu kết quả và đưa bit cộng đầu vào bằng
$s \oplus \rho$: cụ thể $\overline{6} + (1 - 1) = \overline{6} = -7$ (khi $\rho=1$),
còn khi $\rho=0$ thì $\overline{6} + (1 - 0) = -6$. Các trường hợp ngoại lệ
\texttt{Zero}, \texttt{NaR} hoặc kẹp biên $\pm\mathrm{minpos}/\pm\mathrm{maxpos}$
được xử lý song song qua bộ chọn ở ngõ ra.

\paragraph{Hệ quả của vị trí làm tròn đối với MAC-NF, MAC-F và MAC-Quire.}
Đặt $\varepsilon_p = R(\widehat P) - \widehat P$ là sai số làm tròn khi đóng gói
tích về \texttt{Posit(32,3)}, quan hệ giữa giá trị đầu vào bộ cộng của hai kiến
trúc là $R(\widehat P) + C = \widehat P + C + \varepsilon_p$. Khi $\widehat P$ và
$C$ trái dấu và gần triệt tiêu nhau, tổng thực rất nhỏ nên $\varepsilon_p$ trở thành
nhiễu đáng kể làm sai lệch chữ số có nghĩa hoặc ngưỡng kẹp (như đã đo được trên nhóm
$C=R(-P)$ ở Chương~\ref{ch:measurements}). Kiến trúc \texttt{MAC-F} loại bỏ được
$\varepsilon_p$ ở từng phép MAC; tuy nhiên trong chuỗi tích lũy nhiều bước hồi tiếp
ngoài, trạng thái $C$ vẫn phải đi qua một lần lượng tử hóa ở cuối mỗi giao dịch,
trong khi \texttt{MAC-Quire} loại bỏ hoàn toàn bước lượng tử hóa trạng thái trung
gian nhờ thanh ghi $Q$ trên lưới $2^{-480}$.


\section{Lựa chọn kiến trúc MAC-F}
\input{figures/three_architectures}
Hình~\ref{fig:three-architectures} minh họa sự khác biệt về ranh giới làm tròn và
trạng thái giữa ba phương án \texttt{MAC-NF}, \texttt{MAC-F} và \texttt{MAC-Quire}.
Bằng cách đưa trực tiếp tích nội bộ đã chuẩn hóa (dưới dạng magnitude Q24, số mũ
$sf$ và metadata ngoại lệ) vào bộ cộng thay vì phải đi qua một khối Packer trung gian
và Parser giải mã lại như ở \texttt{MAC-NF}, kiến trúc \textbf{\texttt{MAC-F}} loại
bỏ đồng thời hai nguồn tổn hao chính: (i) sai số làm tròn trung gian $R(\widehat P)$
trước khi cộng với $C$, và (ii) chi phí thanh ghi, LUT cùng 3 chu kỳ đường ống phụ
cho khối đóng gói/giải mã tích. Trong khi đó, \texttt{MAC-Quire} loại bỏ hoàn toàn
lượng tử hóa giữa các bước trong chuỗi dài nhưng đánh đổi bằng kích thước ngân
hàng thanh ghi 972 bit, bộ dịch căn lưới rộng và độ trễ giải mã đọc \texttt{READ}
lớn. Dựa trên cân bằng toàn diện giữa chi phí phần cứng, độ trễ và độ chính xác
của từng phép MAC, đồ án lựa chọn \textbf{\texttt{MAC-F}} làm kiến trúc triển khai
và tối ưu sâu, đồng thời sử dụng \texttt{MAC-NF} và \texttt{MAC-Quire} làm mốc
đối chứng định lượng trong Chương~\ref{ch:measurements}.
